\section{竞争、特化与多样性}

\subsection{窄而强 vs 宽而浅}

\begin{importantbox}{超级智能可以是窄的}
Ilya 提出了一个容易被忽视的可能性：超级智能不一定是"什么都会的 oracle"——它可以是\textbf{useful and narrow at the same time}。未来可能有许多不同的窄超级智能 AI，各自专精于不同领域。竞争将通过特化（specialization）展开——就像市场和进化中一样。
\end{importantbox}

\subsection{AI 多样性的问题}

当前 AI 缺乏多样性的原因在于 pre-training——所有模型在相同数据上训练，因此产出惊人地相似。RL 和 post-training 开始引入一些差异化。

\begin{knowledgebox}{Self-Play 与对抗多样性}
Self-play 最初引起 Ilya 兴趣是因为它提供了"仅用 compute、不用 data"生成训练信号的方式。但经典 self-play 只擅长培养特定技能（谈判、冲突、策略）。如今，self-play 以不同的形式回归：debate、verifier、LLM-as-judge 等对抗性设置。

竞争天然催生多样性——当 agent 看到其他 agent 已经采取某种方法时，它有动力去探索不同的方向。
\end{knowledgebox}

\subsection{Timeline 预测}

Ilya 对达到类人持续学习 agent（以及由此而来的超级智能）给出了 \textbf{5--20 年}的时间范围。

\subsection{本章小结}

超级智能可以是窄而强的。竞争将催生特化和多样性。Pre-training 导致了当前模型的同质性，RL 开始引入差异。Ilya 预计 5--20 年达到类人学习能力的 AI。


